\section*{Input-output price shocks and inflation inequality}
\begin{table}[H]
\centering
\caption{Effect of a 40\% oil and gas shock with cross-country propagation, June 2021--June 2023}
\label{tab:shock-oil-gas-chain-ea20}
\scriptsize
\setlength{\tabcolsep}{1.2pt}
\begin{tabular}{lrrrrrr}
\toprule
Country & \shortstack{Effect on average\\inflation -- Total} & \shortstack{Direct effect on\\average inflation} & \shortstack{Indirect effect on\\average inflation} & \shortstack{Effect on\\inflation gap} & \shortstack{Direct effect on\\inflation gap} & \shortstack{Indirect effect on\\inflation gap}
 \\
\midrule
Austria & 11.8 & 3.9 & 7.9 & -1.4 & -0.7 & -0.7
 \\
Belgium & 9.8 & 6.2 & 3.7 & -3.6 & -3.2 & -0.4
 \\
Cyprus & 7.5 & 4.7 & 2.8 & -1.7 & -1.5 & -0.2
 \\
Germany & 7.4 & 4.4 & 3.0 & -1.8 & -1.6 & -0.2
 \\
Estonia & 12.2 & 6.6 & 5.6 & -4.8 & -3.8 & -1.0
 \\
Greece & 6.9 & 4.0 & 2.9 & -2.7 & -2.2 & -0.5
 \\
Spain & 7.4 & 4.5 & 3.0 & -1.9 & -1.6 & -0.3
 \\
Finland & 5.7 & 3.6 & 2.1 & -0.5 & -0.5 & -0.0
 \\
France & 9.3 & 4.5 & 4.7 & -2.8 & -1.9 & -0.9
 \\
Croatia & 11.1 & 7.0 & 4.1 & -3.4 & -2.7 & -0.7
 \\
Ireland & 6.0 & 4.0 & 2.0 & -2.6 & -2.2 & -0.4
 \\
Italy & 9.2 & 4.9 & 4.3 & -1.9 & -1.3 & -0.6
 \\
Lithuania & 8.4 & 6.5 & 1.9 & -2.6 & -2.6 & 0.0
 \\
Luxembourg & 7.7 & 6.0 & 1.8 & -1.1 & -1.1 & -0.1
 \\
Latvia & 14.1 & 7.2 & 6.9 & -7.9 & -4.9 & -3.0
 \\
Malta & 6.1 & 2.9 & 3.2 & -1.0 & -0.7 & -0.3
 \\
Netherlands & 6.4 & 4.3 & 2.0 & -1.4 & -1.3 & -0.0
 \\
Portugal & 8.8 & 3.7 & 5.2 & -4.3 & -2.1 & -2.2
 \\
Slovenia & 8.3 & 5.8 & 2.5 & -1.7 & -1.6 & -0.1
 \\
Slovakia & 12.5 & 6.3 & 6.2 & -0.5 & -0.4 & -0.1
 \\
\midrule
\textbf{Euro area} & \textbf{8.4} & \textbf{4.6} & \textbf{3.8} & \textbf{-2.2} & \textbf{-1.7} & \textbf{-0.5}
 \\
\bottomrule
\end{tabular}
\par\vspace{0.35em}\footnotesize\noindent Sources: Eurostat HICP, HBS, FIGARO, authors' calculations.
\par\vspace{0.25em}\footnotesize\noindent Notes. Entries are percentage-point effects over June 2021--June 2023 of a uniform 40 percent shock across the oil and gas energy chain. The directly shocked FIGARO sectors are B (mining and quarrying, including NACE B06 crude petroleum and natural gas extraction), C19 (coke and refined petroleum products), and D35 (electricity, gas, steam and air-conditioning supply). The shock vector is $s_{c,k}=0.40$ for $k\in\{B,C19,D35\}$ and $s_{c,k}=0$ otherwise, for every country $c$. The price response is $\Delta p = (I - A^{\prime})^{-1}s$, where $A_{(o,i),(d,j)}=Z_{(o,i),(d,j)}/x_{d,j}$ retains bilateral intermediate flows from origin country $o$ and sector $i$ to destination country $d$ and sector $j$. Direct effects arise from household consumption weights mapped to B, C19 and D35; indirect effects include domestic and cross-country propagation through intermediate imports and exports. The euro-area aggregate covers all 20 member countries. Imports from origins outside the EA20 production system enter production costs but are treated as exogenous and do not generate additional endogenous feedback. FIGARO flows are measured in current million euros at basic prices; the dataset used here does not provide separate import or export deflators. Income inflation gap = Q1 - Q5 of revenues.
\end{table}
\vspace{0.6em}
